\paragraph{A nonzero Euclidean response and a clock-sensitive limit test.}
For the same full-Haar coupled source, let
$v=(W_1-\langle W_1\rangle)\Omega$ and
$C(t)=\langle v,e^{-t(H-E_0)}v\rangle$.
All finite spectral weights and a rigorous infinite-representation
remainder give
\[
  0.2355121412297 < C(0)<0.2355121412313,
  \qquad
  0.1077022623741<C(1/4)<0.1079898705277.
\]
The latter interval is a strictly positive full-source response, not a
correlator of a cutoff-defined model.  Three representation resolutions
overlap.  Temporal reflection positivity follows from the positive
source spectral measure, with complete signed reflection forms also
checked in the finite proof coordinates.

At fixed coupling on this fixed graph, the declared physical clock gives
$H(a)=H(1)/a$.  The variance stays positive but
$C_a(t)\leq C(0)e^{-(E_1-E_0)t/a}$ tends to zero for every fixed $t>0$.
This particular sequence fails the noncollapse condition for a strongly
continuous field limit.  A viable four-dimensional Yang--Mills family
must instead supply spatial refinement, controlled coupling flow, common
noncollapsed local fields and the required physical gap scaling, together
with its Euclidean and short-distance identification.  A finite-graph
connected composite correlation does not establish a nonfree continuum.
The separate Lean energy lemma excludes a positive gap only when the
same flat-space vacuum and nonzero finite-energy excitation obey
arbitrary positive dilation covariance; it makes no claim about compact
spatial cylinders or vacua that break that covariance.
